В рамках выпускной квалификационной работы была решена задача реализации и исследования параллельных алгоритмов интерполяции на основе разреженных сеток с квадратичным базисом.

Основным результатом работы стал вычислительный модуль \texttt{compute}, в котором реализованы:
\begin{enumerate}
    \item построение адаптивной разреженной сетки для скалярных и векторных отображений;
    \item поддержка линейного и квадратичного базисов;
    \item последовательный и параллельный режимы построения;
    \item рекурсивное вычисление значения интерполянта;
    \item композиция отображений и итеративный режим через \texttt{MakeNextIteration};
    \item механизм anchor points для предотвращения ложной остановки адаптивного алгоритма;
    \item параметр жёсткого ограничения числа узлов как средство дополнительной оптимизации.
\end{enumerate}

В ходе работы было показано, что разреженные сетки представляют собой эффективный подход к аппроксимации отображений и операторов перехода, возникающих при численном моделировании динамических систем. Отдельно было исследовано применение метода к дифференциальным уравнениям в двух режимах: через включение времени как интервальной неопределённости и через последовательное построение одношагового оператора.

Особое значение имеет предложенный в работе механизм anchor points. На примере функции $\sin(x)$ на интервале $[0,4\pi]$ было показано, что базовый адаптивный критерий может ошибочно интерпретировать функцию как почти линейную, если опорные точки попадают в нули. Реализованный механизм anchor points позволяет в ручном режиме указать важные области и принудить алгоритм к дополнительному уточнению. Это расширение не было заложено в исходной статье \cite{Morozov2024} и является одним из существенных авторских результатов дипломного проекта.

Экспериментальная часть работы показала важность сравнительного анализа различных режимов использования модуля. Были рассмотрены:
\begin{enumerate}
    \item сравнение линейного и квадратичного базисов;
    \item сравнение последовательного и параллельного построения;
    \item влияние размерности задачи и размера области на время построения;
    \item влияние anchor points на корректность и стоимость адаптивного построения;
    \item сравнение двух способов работы с дифференциальными уравнениями;
    \item влияние ограничения числа узлов на вычислительные затраты.
\end{enumerate}

Тем самым дипломная работа привела не только к созданию программной реализации, но и к построению содержательной экспериментальной базы для анализа производительности и границ применимости метода.

Цель работы --- реализация параллельных алгоритмов интерполяции на основе разреженных сеток с квадратичным базисом --- достигнута. Поставленные задачи выполнены в полном объёме.

Практическая значимость работы состоит в том, что созданный вычислительный модуль может использоваться:
\begin{enumerate}
    \item для численного эксперимента с многомерными отображениями;
    \item для исследования операторов перехода динамических систем;
    \item для сравнительного анализа различных стратегий адаптивной интерполяции;
    \item как основа для дальнейшего развития систем визуализации и прикладных сервисов научных вычислений.
\end{enumerate}

Серверный слой на gRPC и клиент визуализации на Flutter расширяют область применения проекта, однако центральным содержанием дипломной работы является именно модуль \texttt{compute}, в котором сосредоточены ключевые алгоритмические результаты.

Перспективы дальнейшего развития проекта связаны с автоматизацией выбора anchor points, дальнейшим исследованием ошибок интерполяции, расширением визуализации на многомерные задачи и построением более развитого контура автоматического анализа результатов бенчмарков.